\part{Interrupts: the axis of everything}

\chapter{What an interrupt is}

\begin{goals}
  \item Why a system without interrupts cannot be an operating system
  \item The three kinds of interrupt and how they differ
  \item What the processor does automatically when one fires
\end{goals}

\section{The problem}

Suppose the kernel wants to know when a key is pressed. Without interrupts there is only one
option: ask, over and over, forever.

\begin{lstlisting}[style=c]
for (;;) {
    if (key_available()) handle_key();
    if (timer_expired()) switch_process();
    if (disk_ready())    finish_read();
    /* ... and the real work? there is no room for it ... */
}
\end{lstlisting}

That loop has two fatal defects. It burns 100\,\% of the processor asking questions whose answer is
almost always ``no''. And it only works if the running code cooperates: if a user program enters
\cod{while(1);}, the loop never runs again and the machine is lost.

\begin{concept}{What an interrupt gives you}
An interrupt lets the hardware \hi{take the processor away} from whatever is running,
without its permission and without its cooperation, and hand it to a kernel routine. That single
capability is what makes preemptive multitasking, device handling, error recovery and system calls
possible. Everything.
\end{concept}

\section{Three kinds}

The processor treats all three the same way, but they originate differently:

\begin{center}
\begin{tabularx}{\textwidth}{@{}l l X@{}}
\toprule
\textbf{Kind} & \textbf{Numbers} & \textbf{Origin} \\
\midrule
\term{Exceptions} & 0--31 & The processor itself, when an instruction cannot complete: divide by zero, invalid opcode, page fault \\
\term{Hardware interrupts (IRQ)} & 32--47 & An external device raising a line: timer, keyboard, disk \\
\term{Software interrupts} & any, we use \hexa{80} & The program itself, executing \cod{int n} on purpose \\
\bottomrule
\end{tabularx}
\end{center}

The difference matters for how you respond. An exception means something went wrong \emph{now}:
either you fix it (a page fault whose page can be brought in) or you kill the culprit. A hardware
interrupt means a device has news. A software interrupt means a program is asking for a service.

\section{What the processor does by itself}

When an interrupt fires, the processor performs a fixed sequence \emph{in hardware}, before a
single instruction of yours runs:

\begin{center}
\begin{tikzpicture}[font=\sffamily\scriptsize,
  p/.style={draw=amber,fill=amberlight,rounded corners=2pt,minimum width=9.4cm,minimum height=0.68cm,align=left,inner xsep=8pt}]
  \node[p] (a) {1. Finish (or abort) the current instruction};
  \node[p,below=0.22cm of a] (b) {2. Look up entry $n$ in the IDT};
  \node[p,below=0.22cm of b] (c) {3. \textbf{If coming from ring 3: switch to the kernel stack from the TSS}};
  \node[p,below=0.22cm of c] (d) {4. Push SS and ESP (only if the ring changed)};
  \node[p,below=0.22cm of d] (e) {5. Push EFLAGS, CS and EIP};
  \node[p,below=0.22cm of e] (f) {6. Push an error code (only some exceptions)};
  \node[p,below=0.22cm of f] (g) {7. Clear IF (with an interrupt gate) and jump to the handler};
  \foreach \x/\y in {a/b,b/c,c/d,d/e,e/f,f/g} \draw[-{Stealth[length=2mm]},amber] (\x)--(\y);
\end{tikzpicture}
\end{center}

Step 3 is the one that makes user space possible, and it is worth dwelling on. If the interrupted
code was running in ring 3, its stack is a user stack --- untrusted, possibly corrupt, possibly
unmapped. The kernel cannot run on it. So the processor automatically switches to a kernel stack
whose address it reads from a structure called the TSS. Part X covers this in detail.

\begin{warn}{Not every exception pushes an error code}
Exceptions 8, 10, 11, 12, 13, 14 and 17 push an extra 32-bit value describing what went wrong. The
others do not. That inconsistency forces the assembly stubs to push a fake zero for the ones that
do not, so that the stack layout is identical in all cases. You will see exactly that trick in
Chapter 21.
\end{warn}

% ============================================================
\chapter{The IDT: the table of handlers}

\begin{goals}
  \item What the Interrupt Descriptor Table is and how it is laid out
  \item How to build a gate descriptor
  \item What the difference between \hexa{8E} and \hexa{EE} is, and why it matters enormously
\end{goals}

\section{256 entries}

The \term{IDT} is an array of 256 descriptors, one per interrupt number. Each entry says which
address to jump to and under what conditions.

\filetag{lytlnybl/include/kernel/interrupts.h}
\begin{lstlisting}[style=c]
typedef struct {
    uint16_t offset_low;    /* bits 15..0 of the handler address */
    uint16_t selector;      /* code segment selector: 0x08 */
    uint8_t  reserved;      /* always 0 */
    uint8_t  flags;         /* type and privileges */
    uint16_t offset_high;   /* bits 31..16 of the handler address */
} __attribute__((packed)) idt_entry_t;
\end{lstlisting}

Note the handler address is split in two halves with other fields in between --- the same
historical mess as the GDT. Filling it in means splitting the address by hand:

\filetag{lytlnybl/kernel/interrupts.c}
\begin{lstlisting}[style=c]
void idt_set_gate(uint8_t interrupt, uint32_t handler_address,
                  uint16_t selector, uint8_t flags) {
    idt[interrupt].offset_low  = handler_address & 0xFFFF;
    idt[interrupt].selector    = selector;
    idt[interrupt].reserved    = 0;
    idt[interrupt].flags       = flags;
    idt[interrupt].offset_high = (handler_address >> 16) & 0xFFFF;
}
\end{lstlisting}

Chapter 2's idioms again: \cod{\& 0xFFFF} to keep the low half, \cod{>> 16} to extract the high
half.

\section{The flags byte}

\begin{center}
\begin{tabularx}{\textwidth}{@{}c l X@{}}
\toprule
\textbf{Bit} & \textbf{Name} & \textbf{Meaning} \\
\midrule
7 & P    & Present \\
6--5 & DPL & \textbf{Minimum ring required to invoke this interrupt with \cod{int}} \\
4 & S    & 0 for gates \\
3--0 & Type & \hexa{E} = 32-bit interrupt gate, \hexa{F} = trap gate \\
\bottomrule
\end{tabularx}
\end{center}

Which produces the two values used in the project:

\begin{center}
\begin{tabularx}{0.95\textwidth}{@{}l l X@{}}
\toprule
\textbf{Value} & \textbf{Binary} & \textbf{Meaning} \\
\midrule
\hexa{8E} & \cod{10001110} & Present, \textbf{DPL 0}, interrupt gate \\
\hexa{EE} & \cod{11101110} & Present, \textbf{DPL 3}, interrupt gate \\
\bottomrule
\end{tabularx}
\end{center}

\begin{concept}{Why the system call gate is \hexa{EE} and everything else \hexa{8E}}
The DPL in an IDT entry answers one question: \emph{may a ring 3 program trigger this interrupt
deliberately with \cod{int n}?}

For the CPU exceptions and the hardware IRQs the answer must be no. If a user program could execute
\cod{int 14} it would fake a page fault and confuse the kernel into examining a \reg{CR2} that has
nothing to do with reality. Hence DPL 0, hence \hexa{8E}.

For \cod{int 0x80} the answer must be yes: it is precisely how a program requests a service. Hence
DPL 3, hence \hexa{EE}.

\hi{That single hex digit is the difference between having system calls and having a
completely inaccessible kernel.} Getting it wrong in either direction is a classic bug: set them
all to \hexa{8E} and every system call becomes a protection fault; set them all to \hexa{EE} and
any program can forge exceptions.
\end{concept}

\filetag{lytlnybl/kernel/interrupts.c}
\begin{lstlisting}[style=c]
idt_set_gate(0,  (uint32_t)isr0,  0x08, 0x8E);   /* divide by zero */
idt_set_gate(14, (uint32_t)isr14, 0x08, 0x8E);   /* page fault */
idt_set_gate(32, (uint32_t)irq0,  0x08, 0x8E);   /* timer */
idt_set_gate(33, (uint32_t)irq1,  0x08, 0x8E);   /* keyboard */
/* ... */
idt_set_gate(0x80, (uint32_t)syscall_entry, 0x08, 0xEE); /* system calls */
\end{lstlisting}

\section{Interrupt gate versus trap gate}

The only difference: an \term{interrupt gate} clears \reg{IF} on entry (interrupts disabled inside
the handler), a \term{trap gate} does not.

lytlnyblOS uses interrupt gates for everything, including system calls. That means \hi{the
kernel is not reentrant}: while a system call is running, no timer interrupt can fire, so no
process switch can preempt it. That simplifies the code enormously --- no locks, no races --- at
the price of latency: a slow system call blocks the whole system.

\begin{warn}{A consequence you will meet later}
Because system calls run with interrupts disabled, the kernel can manipulate process lists,
page directories and the TSS without worrying about being interrupted halfway. Several places in
this codebase rely on that implicitly. Anyone thinking of switching to trap gates would have to
audit all of them first.
\end{warn}

\section{Loading the table}

Like the GDT, the IDT is registered through a small descriptor holding base and limit:

\begin{lstlisting}[style=c]
idtr.limit = sizeof(idt) - 1;    /* 256 * 8 - 1 = 2047 */
idtr.base  = (uint32_t)idt;
idt_load(&idtr);
\end{lstlisting}

\begin{lstlisting}[style=asm]
idt_load:
    mov eax, [esp + 4]
    lidt [eax]
    ret
\end{lstlisting}

The same ``minus one'' as the GDT. And notice \cod{memset(idt, 0, sizeof(idt))} at the start of
\cod{idt\_init}: the 200-odd unused entries are left at zero, which means ``not present''. If any
of them ever fires, the processor raises a general protection fault rather than jumping into
garbage. Zeroing is what makes that safe.

% ============================================================
\chapter{Interrupt service routines}

\begin{goals}
  \item Why handlers cannot be written directly in C
  \item How the assembly stubs build a uniform stack frame
  \item How \cod{registers\_t} maps onto that frame, and why it is the key to everything
\end{goals}

\section{Why C is not enough}

An interrupt handler is not a normal function. It must:

\begin{itemize}
  \item preserve \emph{every} register, because it interrupted arbitrary code at an arbitrary
        point;
  \item end with \cod{iret} instead of \cod{ret}, because the stack holds \reg{EIP}, \reg{CS} and
        \reg{EFLAGS}, not just a return address;
  \item cope with the fact that some exceptions pushed an error code and others did not.
\end{itemize}

C can express none of that. So each of the 48 entry points is a small assembly stub that normalises
the situation and then calls a C function.

\section{The two macros}

\filetag{lytlnybl/kernel/interrupts.asm}
\begin{lstlisting}[style=asm]
%macro ISR_NOERRCODE 1
isr%1:
    push dword 0        ; FAKE error code, so the layout always matches
    push dword %1       ; the interrupt number
    jmp isr_common_stub
%endmacro

%macro ISR_ERRCODE 1
isr%1:
                        ; the processor already pushed the error code
    push dword %1       ; just the interrupt number
    jmp isr_common_stub
%endmacro
\end{lstlisting}

That fake zero is the whole trick. After either macro runs, the stack looks identical regardless of
whether the exception provided an error code --- so a single C function can handle all 32.

\begin{lstlisting}[style=asm]
ISR_NOERRCODE 0     ; divide by zero: no error code
ISR_NOERRCODE 1
...
ISR_ERRCODE   8     ; double fault: has an error code
ISR_NOERRCODE 9
ISR_ERRCODE   10    ; invalid TSS
ISR_ERRCODE   11
ISR_ERRCODE   12
ISR_ERRCODE   13    ; general protection fault
ISR_ERRCODE   14    ; page fault
\end{lstlisting}

\section{The common stub}

\filetag{lytlnybl/kernel/interrupts.asm}
\begin{lstlisting}[style=asm]
isr_common_stub:
    pusha               ; push EAX,ECX,EDX,EBX,ESP,EBP,ESI,EDI in one go

    mov ax, ds
    push eax            ; save the data segment

    push esp            ; the argument: a pointer to everything we just built
    call isr_handler
    add esp, 4          ; discard the argument

    pop eax             ; restore DS (as part of EAX)

    popa                ; restore the eight registers

    add esp, 8          ; discard interrupt number and error code
    iret                ; return: pops EIP, CS, EFLAGS (and SS, ESP if the ring changed)
\end{lstlisting}

\begin{concept}{\cod{push esp} is the most important line in the file}
Right before the call, \reg{ESP} points at the lowest byte of everything the stub and the processor
have pushed. Pushing \reg{ESP} therefore passes to C \hi{a pointer to the complete saved
state of the interrupted code}.

From C, that pointer is a \cod{registers\_t*}. And because it points at the real stack, \emph{it is
not a copy}: if C modifies a field of that structure, it modifies what \cod{popa} and \cod{iret}
will restore. Changing \cod{regs->eip} changes where execution resumes. Changing \cod{regs->eax}
changes the value the interrupted code sees in \reg{EAX}.

This is how system calls return a value. And it is how the context switch works: to switch process,
the kernel simply overwrites every field of that structure with another process's saved values, and
lets \cod{iret} do the rest.
\end{concept}

\section{The \cod{registers\_t} structure}

The structure must match the stack layout exactly, and in reverse order --- the last thing pushed
sits at the lowest address, so it comes first in the struct:

\filetag{lytlnybl/include/kernel/interrupts.h}
\begin{lstlisting}[style=c]
typedef struct {
    uint32_t ds;                /* pushed by the stub */

    uint32_t edi, esi, ebp, esp; /* pushed by pusha */
    uint32_t ebx, edx, ecx, eax;

    uint32_t interrupt_number;  /* pushed by the macro */
    uint32_t error_code;        /* pushed by the macro or by the CPU */

    uint32_t eip;               /* pushed by the processor */
    uint32_t cs;
    uint32_t eflags;

    uint32_t user_esp;          /* only if the ring changed */
    uint32_t ss;
} registers_t;
\end{lstlisting}

\begin{center}
\begin{tikzpicture}[font=\sffamily\scriptsize,x=1cm,y=1cm]
  \def\w{5.6}
  \newcommand{\fila}[4]{%
    \fill[#3,draw=grayborder] (0,#1) rectangle (\w,#1+0.52);
    \node[font=\ttfamily\scriptsize] at (\w/2,#1+0.26) {#2};
    \node[anchor=west,font=\sffamily\tiny,text=gray] at (\w+0.2,#1+0.26) {#4};}

  \fila{0}{ds}{accentlight}{lowest address --- ESP points here}
  \fila{0.52}{edi esi ebp esp}{accentlight}{}
  \fila{1.04}{ebx edx ecx eax}{accentlight}{pushed by \texttt{pusha}}
  \fila{1.56}{interrupt\_number}{amberlight}{pushed by the macro}
  \fila{2.08}{error\_code}{amberlight}{macro or processor}
  \fila{2.60}{eip}{redlight}{}
  \fila{3.12}{cs}{redlight}{pushed by the processor}
  \fila{3.64}{eflags}{redlight}{}
  \fila{4.16}{user\_esp}{greenlight}{}
  \fila{4.68}{ss}{greenlight}{only if ring 3 $\rightarrow$ ring 0}
  \draw[-{Stealth[length=2.5mm]},red,thick] (-1.4,0.26) -- (-0.1,0.26);
  \node[anchor=east,font=\sffamily\scriptsize,text=red] at (-1.5,0.26) {ESP};
\end{tikzpicture}
\end{center}

\begin{warn}{The two last fields are conditional}
\cod{user\_esp} and \cod{ss} only exist on the stack if the interrupt came from ring 3. If the
kernel interrupts itself, the processor does not switch stacks and pushes nothing extra --- so
those fields of the struct point at whatever happened to be above. That is why the context-saving
code checks the process type before touching them:
\begin{lstlisting}[style=c]
if (process->type == PROCESS_KERNEL) {
    process->regs.esp = regs->esp;      /* from pusha */
} else {
    process->regs.esp = regs->user_esp; /* pushed by the processor */
    process->regs.ss  = regs->ss;
}
\end{lstlisting}
Getting this wrong produces a bug that only shows up with user processes, and looks like random
corruption.
\end{warn}

\section{The C handler}

\filetag{lytlnybl/kernel/interrupts.c}
\begin{lstlisting}[style=c]
void isr_handler(registers_t* regs) {
    switch (regs->interrupt_number) {
        case 14:
            page_fault_handler(regs);
            break;
        default:
            vga_text_writeline(&terminal,
                exception_messages[regs->interrupt_number]);
            break;
    }

    if ((regs->cs & 3) == 3) {          /* did it come from user space? */
        current_process->state = PROCESS_TERMINATED;
        context_switch(current_process, get_next_process(), regs);
        return;
    }

    for (;;);                           /* kernel fault: stop dead */
}
\end{lstlisting}

The test \cod{(regs->cs \& 3) == 3} reads the RPL of the saved code segment: if it is 3, the fault
came from a user program. The policy is then clear: kill the program and switch to another. The
system survives.

If the fault came from the kernel itself, there is no safe recovery, so it halts in an infinite
loop. That is deliberate: a frozen system with a message on screen is far easier to debug than one
that keeps running in a corrupt state.

\begin{warn}{The reboot loop this cannot catch}
Notice that the \cod{for(;;)} is only reached \emph{after} printing. If printing itself faults ---
as in the divide-by-zero scenario from Chapter 20 --- the handler never gets there, and the
exception recurses until the stack runs out. A production kernel would use a separate,
pre-allocated fault stack and a panic routine that writes to screen without calling anything. Part
XIII revisits this.
\end{warn}

% ============================================================
\begin{recipe}{adding an interrupt vector of your own}{ch21-custom-interrupt}
Adds vector \hexa{30} to the IDT in the three places it has to go: an assembly stub (C cannot end a
function with \cod{iret}), a declaration so C can take its address, and the \cod{idt\_set\_gate}
call plus a handler.

Then it fires \cod{int \$0x30} with a known value in \reg{EBX}, and the handler overwrites
\cod{regs->ebx}. The value that comes back is the handler's, not the caller's --- demonstrating on
a vector you built yourself the exact mechanism a system call uses to return a result.

Note the handler ends with \cod{return}, not \cod{break}: the code after the switch halts the
kernel forever, which is right for a fault and wrong for a deliberate interrupt.
\end{recipe}

% ============================================================
\chapter{The PIC: routing hardware interrupts}

\begin{goals}
  \item What the 8259 controller does and why there are two of them
  \item Why they must be remapped before anything else
  \item What the End Of Interrupt is and what happens if you forget it
\end{goals}

\section{The chip}

Devices do not signal the processor directly: they raise a line on a \term{Programmable Interrupt
Controller} (PIC), which arbitrates between them and tells the CPU which one won. The original PC
had one, with 8 lines. The PC/AT added a second chained to line 2 of the first, giving 15 usable
lines.

\begin{center}
\begin{tikzpicture}[font=\sffamily\scriptsize,
  d/.style={draw=grayborder,fill=graylight,rounded corners=2pt,minimum width=1.9cm,minimum height=0.5cm},
  p/.style={draw=accent,fill=accentlight,rounded corners=3pt,minimum width=2.1cm,minimum height=2.5cm,align=center}]

  \node[d] (t) at (0,2.1) {Timer (IRQ0)};
  \node[d] (k) at (0,1.4) {Keyboard (IRQ1)};
  \node[d] (s) at (0,0.7) {Slave (IRQ2)};
  \node[d] (c) at (0,0.0) {COM1 (IRQ4)};

  \node[p] (m) at (3.6,1.05) {\textbf{Master}\\PIC 1\\{\tiny 0x20, 0x21}};
  \node[p,minimum height=1.6cm] (sl) at (3.6,-1.9) {\textbf{Slave}\\PIC 2\\{\tiny 0xA0, 0xA1}};
  \node[draw=red,fill=redlight,rounded corners=3pt,minimum width=1.8cm,minimum height=1.1cm] (cpu) at (7.4,1.05) {CPU};

  \foreach \n in {t,k,s,c} \draw[-{Stealth[length=2mm]},gray] (\n) -- (m);
  \draw[-{Stealth[length=2mm]},thick,purple] (sl) -- node[right,font=\tiny,text=gray]{via IRQ2} (m);
  \draw[-{Stealth[length=2.5mm]},thick,red] (m) -- node[above,font=\tiny,text=gray]{INTR} (cpu);
  \node[d,minimum width=1.7cm] at (0,-1.9) {Disk (IRQ14)};
  \draw[-{Stealth[length=2mm]},gray] (0.9,-1.9) -- (sl);
\end{tikzpicture}
\end{center}

\section{Why remapping is mandatory}

By default, the master PIC sends IRQ0--7 as interrupt numbers 8--15. That was fine in 1981. It is a
disaster in protected mode, because \hi{numbers 8--15 are already CPU exceptions}: 8 is
double fault, 13 is general protection fault, 14 is page fault.

Without remapping, a simple timer tick would arrive as ``double fault''. The kernel would print
\cod{Double Fault} a hundred times a second and halt. It is one of the most confusing symptoms a
beginner can hit, because the message is completely unrelated to the cause.

The fix is to reprogram the PICs so they use free numbers: 32--39 for the master, 40--47 for the
slave.

\filetag{lytlnybl/kernel/interrupts.c}
\begin{lstlisting}[style=c]
pic_remap(0x20, 0x28);   /* 32 and 40 */
\end{lstlisting}

\section{The remapping sequence}

The 8259 is configured by sending four \term{initialisation command words} in a fixed order:

\filetag{lytlnybl/kernel/interrupts.c}
\begin{lstlisting}[style=c]
void pic_remap(int offset1, int offset2) {
    uint8_t a1 = inb(PIC1_DATA);       /* save the current masks */
    uint8_t a2 = inb(PIC2_DATA);

    outb(PIC1_COMMAND, ICW1_INIT | ICW1_ICW4);  /* ICW1: start init */
    io_wait();
    outb(PIC2_COMMAND, ICW1_INIT | ICW1_ICW4);
    io_wait();

    outb(PIC1_DATA, offset1);          /* ICW2: vector offsets */
    io_wait();
    outb(PIC2_DATA, offset2);
    io_wait();

    outb(PIC1_DATA, 4);                /* ICW3: slave is on line 2 (bit 2 = 4) */
    io_wait();
    outb(PIC2_DATA, 2);                /* ICW3: slave identity = 2 */
    io_wait();

    outb(PIC1_DATA, ICW4_8086);        /* ICW4: 8086 mode */
    io_wait();
    outb(PIC2_DATA, ICW4_8086);
    io_wait();

    outb(PIC1_DATA, a1);               /* restore the masks */
    outb(PIC2_DATA, a2);
}
\end{lstlisting}

Three details worth noting:

\begin{description}
  \item[The \cod{io\_wait()} between every write] The 8259 is a chip from 1976. Modern processors
    issue writes faster than it can process them. That delay is the trick from Chapter 5.
  \item[ICW3 is asymmetric] The master is told \emph{which line} the slave is on as a bitmask
    ($1 \ll 2 = 4$); the slave is told \emph{its own number} as a plain value (2). Same information,
    two encodings. Pure 1970s design.
  \item[Masks are saved and restored] Each PIC has a register saying which lines are enabled. The
    initialisation sequence would clear it, so the code reads it first and puts it back afterwards.
\end{description}

\section{EOI: the acknowledgement you must not forget}

When a device raises an interrupt, the PIC records that it is \emph{in service} and refuses to
deliver any further interrupt of the same or lower priority until it is told the handler is done.
That acknowledgement is the \term{End Of Interrupt}.

\filetag{lytlnybl/kernel/interrupts.c}
\begin{lstlisting}[style=c]
void pic_send_eoi(uint8_t irq) {
    if (irq >= 8)
        outb(PIC2_COMMAND, PIC_EOI);   /* the slave too, if it came from there */
    outb(PIC1_COMMAND, PIC_EOI);       /* always the master */
}

void irq_handler(registers_t* regs) {
    switch (regs->interrupt_number - 32) {
        case 0: timer_handler(regs); break;
        case 1: keyboard_handler();  break;
        case 2: break;
    }
    pic_send_eoi(regs->interrupt_number - 32);
}
\end{lstlisting}

\begin{danger}{Forgetting the EOI}
It is the classic beginner bug, and its symptom is deeply confusing: \hi{the interrupt fires
exactly once and never again}. The keyboard registers one key and goes dead. The timer ticks once
and multitasking never starts. Nothing crashes, no error appears --- the system simply stops
reacting.

Note also that the EOI must always go to the master, and \emph{additionally} to the slave if the
IRQ came from there. Sending it only to the slave leaves the master permanently blocked.
\end{danger}

\begin{tryit}
With the timer at 100 Hz, adding a counter that prints something every 100 ticks should print once
per second. If it prints once and stops, the EOI is missing or misplaced. If it never prints, the
PIC masks are blocking IRQ0 or \cod{sti} was never executed.
\end{tryit}

\section{Turning interrupts on}

The last line of \cod{idt\_init} is:

\begin{lstlisting}[style=c]
asm volatile("sti");
\end{lstlisting}

\cod{sti} sets \reg{IF}. From this instant the system is interruptible: the timer will start firing,
the keyboard will respond, and --- crucially --- the code that follows can be preempted at any
moment.

Everything before this point ran with interrupts off, in a single-threaded world where nothing could
go wrong concurrently. Everything after it lives in a world where it can.

\begin{summary}
The kernel can now react to the outside world. The next part puts that to use with the two drivers
that make an interactive system possible: the timer, which gives it a heartbeat, and the keyboard,
which gives it ears.
\end{summary}
